\textbf{Теорема Гильберта (о базисе)}

Если $K$ -- нётерово, то $K[x]$ -- нётерово.

\underline{Доказательство.}

Рассмотрим произвольный идеал $I \subset K[x]$. Докажем, что
$I$ конечно порождён (тогда $K[x]$ -- нётерово).

Рассмотрим последовательность многочленов 
(при этом каждый раз будем выбирать многочлены 
\underline{минимально возможной степени}):

$$
f_1 \in I, deg(f_1) = d_1
$$

$$
f_2 \in I \setminus (f_1), deg(f_2) = d_2
$$

$$
\dots
$$

$$
f_n \in I \setminus (f_1, \dots, f_{n - 1}), deg(f_n) = d_n
$$

$$
\dots
$$

Пусть старший коэффициент у $f_1$ -- $a_1$, у $f_2$ -- $a_2$, \dots,
у $f_n$ -- $a_n$, \dots

Предположим, что $K[x]$ не является нётеровым, то есть эта цепочка не
заканчивается.

Рассмотрим цепочку идеалов:

\[
(a_1) \subset (a_1, a_2) \subset \dots \subset
(a_1, a_2, \dots a_n) \subset \dots
\]

Так как $K$ -- нётерово, то эта цепочка стабилизируется, то есть
$\exists N\ :\ a_{N + 1} = b_1 a_1 + b_2 a_2 + \dots + b_N a_N$
($a_{N + 1}$ является линейной комбинацией $a_1, \dots, a_N$).

Теперь посмотрим на многочлен $g$:

$$
g = f_{N + 1} - \sum_{j = 1}^N b_j f_j x^{d_{N + 1} - d_j}
$$

Зачем мы его рассмотрели? Потому что у него коэффициент при $d_{N + 1}$
равен нулю:

$$
a_{N + 1} - \sum_{j = 1}^N b_j a_j = 0
$$

Но $g$ также лежит в $I$, но не лежит в $(f_1, \dots, f_N)$, а значит 
вместо $f_{N + 1}$ мы могли выбрать $g$, у которого степень меньше.
Противоречие с тем, что мы выбирали многочлен $f_{N + 1}$ как
многочлен с наименьшей возможной степенью.
